% introduction.tex -- card v4 scope (eleven test families; P is the only headline family, SAP Addendum A4 Part 2 binds the wording of
% the contribution list). Drafted outcome-neutral on 2026-10-28; on 2026-10-29 the contribution list was completed for the label NEG
% of P (A4.2.4 items 3 and 5). Measured results quoted: the detection table of run P2-GATE-M1 (secondary family S-D, wording of
% decision D0074), the label sentence of P and the clean false-alarm check (run P2-GATE-M4).
% Must-not-claim list (NOVELTY-REAUDIT-GATE section 7; NOVELTY-CHECK-INFFUS section 4, items 10 to 13) is respected: no "first" claim,
% no "decision-level" for the gate, no sentence about Zhou and Lu that only their full text could support.

Input-token gating asks which of the observed tokens given to a model should be withheld before the model reasons from
them.
The question arises for any model that reads a sequence of discrete tokens as evidence, and it is sharp for masked
discrete diffusion models.
A masked diffusion model is trained to predict any masked token from any visible subset of the sequence, and the mask
token means ``unknown''~\cite{austin2021d3pm,sahoo2024mdlm,ou2025radd}.
Every token that is visible is therefore treated as certain.
This is acceptable for a clean training sequence.
It is a weakness when the input comes from a pipeline that compresses, adds noise to, crops, splices or mislabels the data,
because a single wrong visible token is used with the same confidence as a correct one.
Unified models extend the setting to two or more modalities in one token sequence~\cite{swerdlow2025unidisc,yang2025mmada,li2025lavida,shi2025muddit}.
In such a model a damaged image token can change what the model concludes about a caption, and the reverse.

\paragraph{Why this is a fusion problem.}
Each image token and each caption token is a source that reports a value, and the model combines them.
The paired decision studied here, whether a caption describes an image, sums per-token log-likelihood ratios over the 1\,024
image tokens of one image.
That is the rule used to fuse local decisions in distributed detection, in which each source is weighted by its reliability
before a global threshold is applied~\cite{chair1986optimal,tenney1981detection}; it is optimal for sources that are independent
given the hypothesis, which image tokens are not, so here it is an approximation (Section~\ref{ssec:limits}).
Imperfect and spurious data are a first-order difficulty of multisensor fusion~\cite{khaleghi2013multisensor,vempaty2013byzantine},
a survey in this journal organises multimodal fusion on low-quality data into noisy, incomplete, imbalanced and quality-varying
settings~\cite{zhang2026lowquality}, and evidential fusion methods discount each modality by an estimated
reliability~\cite{shao2024ddef,huang2025evidential}.
Classical practice also checks each source for a fault and isolates the faulty ones before fusing~\cite{willsky1976survey,gao2015survey},
and sets the detection threshold from clean data so that the false-alarm rate stays at a chosen value, as in constant false-alarm
rate (CFAR) detection~\cite{finn1968adaptive,rohling1983cfar}.
Here the image tokens are the sources, the masked denoiser is the model of what a healthy source looks like, and the gate is the
fault detector that runs before the evidence is summed.
We call this input-level source validation: the gate acts on the sources before the evidence sum, not on the output decision.
The paper proposes no new fusion rule; it specifies a gate, a threshold with a stated guarantee, and a preregistered measurement
of whether withholding the flagged tokens changes the decision.

\paragraph{A measured failure case.}
The gate helps only if the model's own masked likelihood separates damaged tokens from intact ones, and a detection measurement
made before the main test shows that it does so for one kind of damage and not for another.
% run: P2-GATE-M1
Tokens replaced by uniform random codes are separated almost perfectly (mean per-item AUROC 0.961 for a replaced patch and 0.955
for swapped codes), whereas the score is close to chance on JPEG compression (0.561) and Gaussian noise (0.503), where the changed
codes depend on the image content.
This is a measurement of the score in a secondary family (S-D, Section~\ref{ssec:res_det}); it does not enter the headline test and
is not used to interpret it.

\paragraph{What exists and what is missing.}
Many recent methods remask tokens in a masked diffusion model.
Most act on tokens that the sampler itself produced: they remask low-confidence commitments, learn a quality score, or
perturb the context to find brittle tokens
\cite{wang2025remdm,kim2025prism,schiff2026proseco,huang2025remedi,zhai2026core,yao2026remask,lezama2022tokencritic}.
In multimodal masked diffusion models the same idea has been applied to the model's own generations, with cross-modal
confidence as the trigger~\cite{ouyang2026freecorrection,le2026co2jump}.
Work on corrupted \emph{input} tokens is thinner.
Trust but Verify~\cite{jhingran2026trustverify} detects corrupted conditioning tokens in a text-only masked diffusion model,
re-noises them to the mask state and refills them; it needs a finetune, uses a hand-set threshold, and has no second modality.
Conformal risk control has been used with masking-derived scores to flag erroneous caption words given an image, with a
CLIP-based score and a threshold calibrated on labelled foil data~\cite{gomes2025crc}.
In this journal, Zhou and Lu~\cite{zhou2027burden} describe, in their abstract, a keep-or-override gate between two
vision-language recognisers whose threshold is chosen with a Clopper--Pearson bound on harmful overwrites; its unit is a whole
prediction rather than an input token.
The authors of UniDisc show repair of corrupted image-caption pairs by random masking and best-of-$n$ re-infilling in one
qualitative figure~\cite{swerdlow2025unidisc}.
Section~\ref{sec:related} gives the precise differences.
Three things are missing from this picture.
No measurement separates a gate that finds damaged tokens from a mask that merely has the same shape.
No threshold for input tokens carries a stated bound on the false-alarm proportion on clean inputs.
And no comparison shows whether withholding flagged tokens helps a cross-modal decision more than cheap non-generative
detectors or a pixel denoiser do.

\paragraph{Contributions, stated as a measurement.}
This paper does not claim to be the first to detect or remask corrupted input tokens, the first to remask in a unified
model, the first risk-controlled threshold for vision-language errors, or the first masked-likelihood or null-normalised
generative score.
The gate is the object of measurement here, not an asset whose merit is assumed.
Each numbered item corresponds to a result that the preregistered analysis delivered whatever it found.
\begin{enumerate}
\item \emph{The measured object} (Section~\ref{sec:method}). We specify an input-level source-validation gate for a pretrained
unified masked diffusion model: each image token is scored by the model's own masked likelihood with the caption withheld, and
the flag threshold is set by conformal risk control on clean pairs so that the expected flagged share of a clean pair's tokens is
at most $\alpha$.
\item \emph{A closed-form condition} (Section~\ref{ssec:evidence}). Under an independent-token approximation, we give the condition
on the false-alarm and hit rates under which discarding flagged tokens raises the signal-to-noise ratio of an additive evidence
sum; it is an approximation because image tokens are dependent, and it is not linked to any result of the primary test.
\item \emph{A preregistered test} (Sections~\ref{sec:protocol} and~\ref{ssec:res_primary}). An exact McNemar test with Holm
correction over three realistic corruptions, with the sample size fixed by a blind headroom rule, of whether withholding the
flagged tokens changes image-direction verification accuracy, with controls for mask geometry, texture and pixel denoising.
% A4.2.4 item 3, label NEG: the first clause of the opening of Section 6.1 (A4.2.3). run: P2-GATE-M4 (outcome.label = negative)
The preregistered test excludes an improvement or a loss of more than five percentage points in image-direction verification accuracy
from gating under JPEG quality~10, Gaussian noise with $\sigma=0.10$ and a $128$-pixel splice.
\item \emph{Token detection, measured} (Section~\ref{ssec:res_det}).
% run: P2-GATE-M1
At the token level the gate's score separates code-level random replacements (AUROC 0.961 for a patch, 0.955 for swapped codes)
but is near chance on JPEG and Gaussian noise (0.561 and 0.503).
On noise it is below a one-pass score without finetuning (0.607; the marginal 95\,\% intervals, 0.493--0.513 and 0.597--0.618, do
not overlap; no paired test was run), and on JPEG the two overlap (0.561, 0.547--0.575, against 0.581, 0.569--0.593).
In the preregistered comparison it is below a texture score whose sign was chosen on the same items (0.756 and 0.768), so the
preregistered falsifier for the score was met.
\item \emph{Calibration, checked} (Section~\ref{ssec:res_cal}). We check the clean false-alarm statement on held-out clean test
pairs.
% run: P2-GATE-M4 (calibration.fap_check.per_alpha.0.05)
At $\alpha=0.05$ the realised flagged share was 0.0456 (95\,\% interval 0.0442 to 0.0470), with no evidence of violation.
% A4.2.4 item 6: list only those of S2 and B whose status is "run, formal" or "run, descriptive", with that status; omit the item if
% neither ran. R was not run (gatekept: no primary rejection with b>c), so it is not listed.
\item \todo{RESULT (A4.2.4 item 6): \emph{Secondary families} (Sections~\ref{ssec:res_s2} and~\ref{ssec:res_text}): name only those
of re-infilling (S2) and two-way fusion on hard negatives (B) that ran, with their status; delete this item if neither ran.}
\end{enumerate}
The scope is one pretrained 1.4B model, one image-caption dataset, three realistic image corruptions in the primary family and
two sanity corruptions in a separate positive control (PS).
The primary endpoint is the image-direction statistic, which has headroom on clean pairs; the text direction is near ceiling
there, and Section~\ref{ssec:task} states what this choice limits.
Caption corruption is not tested.

\todo{AUTHOR: Figure 1 (motivation, page 1-2 per IF anatomy rule 31) is being prepared separately from committed run arrays; do not
draw it with a generative image tool.}

\paragraph{Roadmap.}
Section~\ref{sec:related} places the work against the closest methods.
Section~\ref{sec:method} defines the scores, the calibrated gate, the evidence statistic and the algorithms.
Section~\ref{sec:protocol} gives the data, folds, corruptions, arms, controls, the families and the reading of the primary test;
the preregistration record and the deviations are in \ref{app:prereg}.
Section~\ref{sec:results} reports the results family by family, Section~\ref{sec:discussion} discusses them with the limitations,
and Section~\ref{sec:conclusion} concludes.
\ref{app:timefree} gives the evidence that the model has no time input in the released configuration.
